\section{Mejora de una estimación matricial del Lema 7.5 de C}
Sea $P$ bistocástica con $\max P_{ij}\le M_0/K$, y $B=P-J_K$ con
$\|B\|_{\rm op}\le r<1$, donde $J_K$ tiene entradas $1/K$. Estos son los parámetros
matriciales del último tablero de C. Sean
\[
 S_0=\begin{pmatrix}I&P\\P^T&I\end{pmatrix},\quad
 T=\begin{pmatrix}0&B\\B^T&0\end{pmatrix},\qquad
 \Delta=\sum_e\frac{P_e^2}{1-2P_e}v_e v_e^T,
\]
donde $v_e$ indica las dos líneas RC de la casilla. Como $T$ anula las dos direcciones
constantes, con $\Pi_\pm$ las proyecciones sobre $(\mathbf1,\pm\mathbf1)$,
\[
 S_0^+=(I+T)^{-1}-\tfrac12\Pi_+-\Pi_- .
\]
La identidad $(I+T)^{-1}-I=-T+T(I+T)^{-1}T$ y la bistocasticidad dan
$\|Te_j\|_2^2\le (M_0-1)/K$ para todo $j$. Como
$\|(I+T)^{-1}\|\le(1-r)^{-1}$, cada entrada de $S_0^+-I$ es $O(1/K)$.
Todos los $v_e$ tienen norma $\ell^1$ dos. Para $K> 2M_0$,
$\sum_e P_e^2/(1-2P_e)\le M_0/(1-2M_0/K)$. Por tanto
\[
 |\tr((S_0^+-I)\Delta)|=O_{M_0,r}(K^{-1}).
\]
El resto cuadrático de $\log\pdet(S_0+\Delta)$, el desarrollo del factor local y las
otras trazas de la prueba del Lema 7.5 de C ya tienen resto $O(1/K)$.
Esta mejora elimina la pérdida $K^{-1/2}$ de la comparación del normalizador con
la expresión determinantal explícita. El error del permanente de la Proposición 7.4
es una entrada distinta. La mejora por sí sola no implica un desarrollo de $\log\per$
con resto $O(K^{-2})$.
